\section{Tái lập}
\label{sec:tai-lap}

Đề yêu cầu: cố định seed; báo cáo trung bình $\pm$ độ lệch chuẩn trên ít nhất 5 lần
chạy độc lập hoặc 5 fold; siêu tham số chỉ chọn trên dữ liệu kiểm định
\cite[tr.~2, mục II]{de}. Mục này ghi điều kiện tái lập và nêu rõ phần nào còn
thiếu bằng chứng.

\subsection{Cấu trúc kho mã}

\begin{center}
\begin{tabular}{lp{0.60\textwidth}}
\toprule
\textbf{Đường dẫn} & \textbf{Nội dung}\\
\midrule
\texttt{src/} & Mã dùng lại: \texttt{attention.py} (A.1), \texttt{char\_lm.py}
(A.2), \texttt{lora.py} (A.3), \texttt{dpo.py} (A.4), \texttt{d3\_data.py}\\
\texttt{notebooks/} & Notebook điều khiển và giải thích, chạy tuần tự trong phiên
sạch\\
\texttt{tests/} & Bộ kiểm tự động\\
\texttt{scripts/} & Chuẩn bị dữ liệu D.2, đóng gói và huấn luyện D.3\\
\texttt{docs/theory/} & Phần B biên dịch độc lập\\
\texttt{docs/report/} & Báo cáo này\\
\texttt{docs/data/} & Nhật ký nguồn, quyền, quy trình riêng tư, datasheet\\
\texttt{docs/progress/} & Ma trận yêu cầu, bối cảnh, kế hoạch, nhật ký AI\\
\texttt{references/} & Năm bài báo và PDF đề gốc\\
\bottomrule
\end{tabular}
\end{center}

\subsection{Môi trường}

\textbf{Môi trường CPU cục bộ} (dùng cho toàn bộ phép kiểm A.1, A.3, A.4 và C.1--C.4):

\begin{center}
\begin{tabular}{ll}
\toprule
Python & 3.14.3\\
PyTorch & 2.12.1+cpu\\
CUDA khả dụng & \texttt{False}\\
\bottomrule
\end{tabular}
\end{center}

Bản PyTorch cài cục bộ là bản CPU, nên mọi con số ở các mục trên là số CPU. Máy có
GPU rời nhưng bản PyTorch hiện tại không dùng được nó; nhóm \emph{không} cài lại
PyTorch tùy tiện vì điều đó làm mất khả năng tái lập các kết quả đã ghi.

\textbf{Môi trường GPU cho D.3} khai báo ở \texttt{requirements-d3.txt}:
\texttt{transformers==4.57.6}, \texttt{peft==0.18.0},
\texttt{accelerate==1.12.0}, \texttt{bitsandbytes==0.48.2}. PyTorch và CUDA lấy từ
image Kaggle và \emph{không} thay thế, để tránh lệch giữa bản đã kiểm và bản chạy.

\textbf{Môi trường A.2} là một phiên Kaggle. Phiên hiển thị cấu hình
\texttt{GPU T4 x2} và trạng thái \emph{Version \#1 --- Successful}.

\subsection{Bộ kiểm tự động \dalam}

Chạy từ gốc kho mã:

\begin{quote}\small\ttfamily
python -m pytest -q
\end{quote}

Kết quả lần chạy ghi trong báo cáo này: \textbf{30 test qua, 1 test bỏ qua, 13
subtest qua} trong $8{,}38$ giây trên môi trường CPU nêu trên. Test bị bỏ qua là
test cần phụ thuộc không có trong môi trường cục bộ.

\subsection{Seed và số lần chạy}

\begin{center}
\begin{tabular}{lll}
\toprule
\textbf{Hạng mục} & \textbf{Seed} & \textbf{Số lần lặp}\\
\midrule
A.2 huấn luyện & 41, 42, 43, 44, 45 & 5 lần huấn luyện độc lập\\
C.1 đầu vào minh họa & \texttt{577893521} & 1 (bộ xác định, không lấy mẫu lại)\\
C.4 đầu vào minh họa & --- & 3 dấu margin $\times$ 4 giá trị $\beta$ (xác định)\\
\bottomrule
\end{tabular}
\end{center}

\textbf{Đơn vị lặp của A.2} là \emph{một lần huấn luyện hoàn chỉnh}, không phải một
lần sinh từ cùng checkpoint. Độ lệch chuẩn báo cáo là SD mẫu (mẫu số $n-1=4$) trên
5 giá trị PPL. Đây là độ biến thiên do khởi tạo và thứ tự trộn \emph{trên cùng một
split}; nó không phải khoảng tin cậy và không phải đánh giá trên tập test độc lập.

\subsection{Tách dữ liệu và quy tắc chọn siêu tham số}

\begin{itemize}[leftmargin=2em]
\item \textbf{A.2:} tách train (prefix 90\%) và validation (suffix 10\%) theo code
point \emph{trước} khi học từ vựng và \emph{trước} khi tạo cửa sổ. Từ vựng chỉ học
từ train. Phần đuôi không đủ một cửa sổ bị loại.
\item \textbf{D.2:} tách train/validation/test, kèm tách nhóm gần trùng (cùng họ
thủ tục, cùng câu hỏi). Tập test là \texttt{USER\_TEST}, có công cụ kiểm riêng để
chặn rò rỉ sang train.
\item \textbf{Quy tắc chung:} không dùng tập test để chọn cấu hình. Siêu tham số
của D.3 sẽ chọn trên validation; kế hoạch so $r$, $\alpha$, LR, epoch đã viết theo
quy tắc đó.
\end{itemize}

\subsection{Những gì còn thiếu để tái lập đầy đủ \chualam}

Nhóm liệt kê thẳng, thay vì tuyên bố đã đạt:

\begin{enumerate}[leftmargin=2em]
\item \textbf{Log môi trường và checkpoint của A.2.} Tệp \texttt{summary.json}
nhận được không chứa phiên bản PyTorch/CUDA và không kèm checkpoint. Nhóm chưa đối
chiếu được số GPU thực dùng. Việc cần làm: lưu log và output của phiên Kaggle đó.
\item \textbf{Hiện vật pilot D.2 không có trên máy hiện tại.} Các bản lưu
(\texttt{data/review/}, \texttt{.cache/}) nằm ngoài Git theo quy tắc riêng tư, nên
chỉ còn báo cáo JSON đã commit. Người tái lập cần nhận bàn giao riêng.
\item \textbf{Toàn bộ số đo của D.3--D.5} chưa tồn tại, nên chưa có gì để tái lập.
\item \textbf{Yêu cầu 5 lần chạy chưa áp dụng cho Phần D.} Kết quả 5 seed hiện chỉ
thuộc A.2. Không được suy ra rằng yêu cầu này đã đạt cho phần D.
\end{enumerate}

\subsection{Cách chạy lại từng phần}

\begin{center}
\begin{tabular}{ll}
\toprule
\textbf{Mục} & \textbf{Lệnh hoặc notebook}\\
\midrule
A.1 & \texttt{notebooks/01\_attention\_decoder.ipynb}\\
A.4 & \texttt{notebooks/02\_dpo\_loss.ipynb}\\
A.3 & \texttt{notebooks/03\_lora\_linear.ipynb}\\
C.2 & \texttt{notebooks/05\_transformer\_block\_cost.ipynb}\\
C.3 & \texttt{notebooks/06\_lora\_qlora\_cost.ipynb}\\
C.4 & \texttt{notebooks/07\_dpo\_loss\_gradient.ipynb}\\
C.1 & \texttt{notebooks/08\_causal\_attention\_example.ipynb}\\
D.3 & \texttt{notebooks/09\_d3\_sft\_kaggle.ipynb} (cần release D.2)\\
Phần B & \texttt{tectonic docs/theory/part-b.tex}\\
Báo cáo & \texttt{tectonic docs/report/report.tex}\\
\bottomrule
\end{tabular}
\end{center}

Các notebook C.1--C.4 không cần GPU, mạng, trọng số hay dữ liệu QA; chúng chạy tuần
tự trên CPU và tự tìm gốc kho mã.
